\section*{Problema 20 del Capítulo 4}

\textbf{Enunciado:}

Sea
\[
    f(x) = \begin{cases} 
        0, & \text{si } x \text{ es irracional}, \\
        \frac{1}{q}, & \text{si } x = \frac{p}{q} \text{ es racional en forma irreducible.}
    \end{cases}
\]

(Un número $p/q$ es \textit{irreducible} si $p$ y $q$ son enteros sin divisores comunes, y $q > 0$). Dibujar la gráfica de $f$ tan bien como se sepa (no desparramar puntos al azar sobre el papel; consídérese en primer lugar los números racionales con $q = 2$ y después aquéllos con $q = 3$, etc.).

\textbf{Gráfica:}

\begin{figure}[h!]
    \centering
    \begin{tikzpicture}
        \begin{axis}[
            axis lines = middle,
            xlabel = $x$,
            ylabel = {$f(x)$},
            ymin = 0, ymax = 1.1,
            xmin = 0, xmax = 1,
            width=12cm,
            height=8cm
        ]
        % Gráfica de los puntos racionales con q=2
        \addplot[only marks, mark=*, mark size=1pt] coordinates {
            (0.5, 1/2)
        };
        % Gráfica de los puntos racionales con q=3
        \addplot[only marks, mark=*, mark size=1pt] coordinates {
            (1/3, 1/3) (2/3, 1/3)
        };
        % Gráfica de los puntos racionales con q=4
        \addplot[only marks, mark=*, mark size=1pt] coordinates {
            (1/4, 1/4) (3/4, 1/4)
        };
        % Gráfica de los puntos racionales con q=5
        \addplot[only marks, mark=*, mark size=1pt] coordinates {
            (1/5, 1/5) (2/5, 1/5) (3/5, 1/5) (4/5, 1/5)
        };
        % Gráfica de los puntos racionales con q=6
        \addplot[only marks, mark=*, mark size=1pt] coordinates {
            (1/6, 1/6) (5/6, 1/6)
        };
    \end{axis}
    \caption{Gráfica de $f(x)$ para los números racionales en forma irreducible.}
    \label{fig:funcion_f}
\end{figure}

\end{document}